\section{Target User Research}

\subsection{Age Ranges and Skill Levels of Target Users}

The DSL's target users are early adolescents, approximately 10--15 years of age, who are mature enough to read and type but are still novices in programming. 
The age ranges were inspired from studying existing programs; for example, the new Irish curricula introduce coding in ``Junior Cycle'' (approx. 12--15 years of age) \cite{noone_visual_2018}. 
These users have usually mastered basic computational thinking (often via visual tools) and are ready to tackle syntax-based tasks. 
The DSL being text-based implies that the learners must have sufficient typing skills, although it assumes no prior familiarity with programming syntax, as it is designed for learning. 
The DSL is aimed at primary-school and middle-school aged children transitioning from block-based environments toward text, who can type steadily and are building their first independent coding skills \cite{noone_visual_2018, saito_comparison_2017}.

\subsection{Accessibility Requirements and Learning Challenges}

Accessible design must accommodate young learners. 
From a learning standpoint, novices universally face conceptual hurdles. 
Reviews find that beginners often lack familiarity with syntax, algorithms, and debugging strategies, as they struggle with syntactic, conceptual, and strategic knowledge in introductory programming \cite{neuwinger_systematic_2025}. 
Common errors include forgetting punctuation or misunderstanding loops. 
The DSL must thus minimize unnecessary complexity through clear error messages and gradual feature introduction, and build on children's math and problem-solving skills, since those are known challenges in learning to code \cite{neuwinger_systematic_2025, nadeem_enhancing_2021}.

\subsection{Influence of the DSL on Adjacent Users}

The choice of a text-based DSL affects not only the target user but also those around them. 
Text-based code influences novices to rely on the teacher for syntax help, so instructors become more involved in technical guidance \cite{weintrop_comparing_2018}. 
The DSL shapes classroom culture; an intuitive, playful coding environment makes the class friendly and encourages peer interaction \cite{weintrop_comparing_2018}. 
In practice, this means that if the DSL is engaging, siblings or classmates might be more inclined to join in or learn together. 
Likewise, parents may take on roles as helpers or reviewers of the child's code, and their experience will be influenced by how approachable the text language is. 
In sum, the DSL influences educators' teaching style and peers' collaboration.

\subsection{Stakeholders Involved}

Besides the children, multiple stakeholders are involved in the solution. 
Key educational stakeholders include teachers, school administrators, and parents \cite{akkaya_examining_2025}, since they guide and support the child's learning process. 
Curriculum designers and educational researchers are also stakeholders, as they determine how a new language fits into learning standards and lesson plans \cite{noauthor_stakeholders_2022}. 
In practice, administrators and policymakers decide whether programming is part of the school curriculum, directly affecting the DSL's adoption. 
In summary, the ecosystem includes the student coders, their families, their teachers, the school leaders, and the broader education community --- curriculum experts, policymakers, and developers --- all of whom play roles in the design, deployment, and success of the DSL \cite{akkaya_examining_2025, noauthor_stakeholders_2022}.
